\section{Modeling Three Dimensions in Lower Dimensions}

The selection argument so far reads as if the lower-dimensional substrates were impossible. They are not. Their
missing features could be made to emerge at a higher layer, and a complete account has to say so. This caveat does
not overturn the selection. It sharpens it. The right reading of the previous chapters is not that the
alternatives are forbidden, but that $\{3,4,3,4\}$ is the substrate of \emph{least emergence}. This chapter makes
that distinction precise.

\subsection{Can three-dimensional physics live in a lower base}

It can, in three different senses, and none of them is forbidden by any principle the model relies on.

\begin{itemize}
  \item Computation. A one-dimensional cellular automaton can be Turing-complete, as Rule 110 is, so it can
  simulate three-dimensional physics outright. In that case the three-dimensionality lives in the data, not in the
  geometry of the base.
  \item Holography. A gravitational theory in $d+1$ dimensions is exactly equivalent to a theory on a
  $d$-dimensional boundary. The $\mathrm{AdS}_3/\mathrm{CFT}_2$ correspondence, for instance, puts a
  three-dimensional gravity inside a two-dimensional conformal field theory. The extra dimension emerges from the
  boundary data.
  \item Emergent dimension. Tensor networks such as MERA, the SYK model, and matrix models all grow effective
  dimensions out of lower-dimensional bases.
\end{itemize}

The reason any of this is possible is the holographic bound. Information scales with boundary area, not with
volume, so three-dimensional physics can in principle be fit onto a lower-dimensional screen. Dimension is not a
hard wall.

\subsection{The catch}

The escape is real but it is not free. A naively local lower-dimensional theory cannot do it. The Bell
inequalities and the Weinberg-Witten theorem both block a straightforward local encoding, so the encoding has to
be non-local, or it has to grow an emergent dimension. In practice a lower-dimensional substrate that wanted to
reach three-dimensional fermionic physics would have to carry extra machinery, a holographic or internal extra
dimension, a fermionization to manufacture the spin statistics, and an emergent gauge structure. Each of these is
a layer that the base does not have on its own.

\subsection{Why the four-dimensional substrate still wins}

This is where the economy argument lands. $\{3,4,3,4\}$ hands over three-dimensional space, spin, and gauge at the
base, with no dimensional emergence required at all. The lower-dimensional substrates could in principle reach the
same physics, but only by paying an emergence cost, uplifting the dimension, fermionizing, and generating a gauge
group. The selection is therefore made on economy rather than on impossibility.

\begin{principle}[Least emergence]
The substrate is selected by economy, not by exclusion. $\{3,4,3,4\}$ delivers three-dimensional space, spin, and
gauge at the base with no emergent dimension. A lower-dimensional substrate could reach the same physics only by
adding emergent machinery. The precise claim is that $\{3,4,3,4\}$ is the least-emergence substrate, not the only
conceivable one.
\end{principle}


Stating the argument this way keeps it defensible and keeps it falsifiable. The position is not "the others are
impossible", which would be hard to defend and easy to attack. It is "the others need more emergent layers to do
what $\{3,4,3,4\}$ does at the base", which is a measured comparison and an open invitation to check. It also
keeps the door open in two useful directions. The emergence question stays alive as a research target, and the
lower-dimensional substrates keep their role as controls. Whether the crystal is finite or infinite is a separate
matter, and it is left open here as elsewhere in the paper.
